%% ++++++++++++++++++++++++++++++++++++++++++++++++++++++++++++
%% Hauptdatei, Wurzel des Dokuments
%% ++++++++++++++++++++++++++++++++++++++++++++++++++++++++++++

% Headerfeld, Typ des Dokumentes, einzubindende Packages.
% Hier bei Bedarf Änderungen vornehmen.
\documentclass
[   twoside=false,     % Einseitiger oder zweiseitiger Druck?
    fontsize=12pt,     % Bezug: 12-Punkt Schriftgröße
    DIV=15,            % Randaufteilung, siehe Dokumentation "KOMA"-Script
    BCOR=17mm,         % Bindekorrektur: Innen 17mm Platz lassen. Copyshop-getestet.
%    headsepline,
    headsepline,  % Unter Kopfzeile Trennlinie (aus: headnosepline)
    footsepline,  % Über Fußzeile Trennlinie (aus: footnosepline)
    open=right,        % Neue Kapitel im zweiseitigen Druck rechts beginnen lassen
    paper=a4,          % Seitenformat A4
    abstract=true,     % Abstract einbinden
    listof=totoc,      % Div. Verzeichnisse ins Inhaltsverzeichnis aufnehmen
    bibliography=totoc,% Literaturverzeichnis ins Inhaltsverzeichnis aufnehmen
    titlepage,         % Titelseite aktivieren
    headinclude=true,  % Seiten-Head in die Satzspiegelberechnung mit einbeziehen
    footinclude=false, % Seiten-Foot nicht in die Satzspiegelberechnung mit einbeziehen
    numbers=noenddot   % Gliederungsnummern ohne abschließenden Punkt darstellen
]   {scrreprt}         % Dokumentenstil: "Report" aus dem KOMA-Skript-Paket

\usepackage[active]{srcltx}
%\usepackage[activate=normal]{pdfcprot} % Optischer Randausgleich -> pdflatex!
\usepackage{ifthen}
\usepackage[ngerman]{babel}   % Neue Deutsche Rechtschreibung
%\usepackage[latin1]{inputenc} % Zeichencodierung nach ISO-8859-1
\usepackage[utf8]{inputenc}   %	Zeichencodierung nach UTF-8 (Unicode)
\usepackage[T1]{fontenc}
%\usepackage{ae} % obsolet und durch lmodern ersetzt
\usepackage{lmodern}
\usepackage[T1]{url}
\usepackage[final]{graphicx}
\RequirePackage{scrlfile}
\ReplacePackage{scrpage2}{scrlayer-scrpage}
% old: \usepackage[automark]{scrpage2}
\usepackage[automark]{scrlayer-scrpage}
\usepackage{setspace}
%\usepackage[first,light]{draftcopy} % Für Probedruck
\usepackage[plainpages=false,pdfpagelabels,hypertexnames=false]{hyperref}

% Tiefe der Kapitelnummerierung beeinflussen
\setcounter{secnumdepth}{3} % Tiefe der Nummerierung
\setcounter{tocdepth}{3}    % Tiefe des Inhaltsverzeichnisses

% Datum anpassen
\newcommand{\leadingzero}[1]{\ifnum #1<10 0\the#1\else\the#1\fi}    
\renewcommand{\today}{\leadingzero{\day}.\leadingzero{\month}.\the\year}     % DD.MM.YYYY

% Hier in die zweite geschweifte Klammer jeweils
% die persönlichen Daten und das Thema der Arbeit eintragen:
\newcommand{\artderausarbeitung}{C - Trainer Arbeit}
\newcommand{\namedesautors}{Jonathan Skopp}
\newcommand{\themaderarbeit}{Starke und schwache Felder}

% PDF Metadaten definieren
\hypersetup{
   pdftitle={\themaderarbeit},
   pdfsubject={\artderausarbeitung},
   pdfauthor={\namedesautors},
   pdfkeywords={\artderausarbeitung; TU-Ilmenau}}


% Abkürzungsverzeichnis beeinflussen. Hier nichts ändern!
\usepackage[intoc]{nomencl}
  \AtBeginDocument{\setlength{\nomlabelwidth}{.25\columnwidth}}
  \let\abbrev\nomenclature
  \renewcommand{\nomname}{Abkürzungsverzeichnis und Formelzeichen}
  \renewcommand{\nomlabel}[1]{#1 \dotfill}
  \setlength{\nomitemsep}{-\parsep}
  \makenomenclature

\usepackage[normalem]{ulem}
  \newcommand{\markup}[1]{\textbf{#1}}

% Seitenlayout festlegen. Hier nichts ändern!
\pagestyle{scrplain}
\ihead[]{\headmark}
\ohead[]{\pagemark}
\chead[]{}
\ifoot[]{}
\ofoot[]{\scriptsize \artderausarbeitung\ \namedesautors}
\cfoot[]{}
\renewcommand{\titlepagestyle}{scrheadings}
\renewcommand{\partpagestyle}{scrheadings}
\renewcommand{\chapterpagestyle}{scrheadings}
\renewcommand{\indexpagestyle}{scrheadings}

% Abschnittsweise Nummerierung anstatt fortlaufend. Hier nichts ändern!
\makeatletter
\@addtoreset{equation}{chapter}
\@addtoreset{figure}{chapter}
\@addtoreset{table}{chapter}
\renewcommand\theequation{\thechapter.\@arabic\c@equation}
\renewcommand\thefigure{\thechapter.\@arabic\c@figure}
\renewcommand\thetable{\thechapter.\@arabic\c@table}
\makeatother

% Quelltextrahmen, klein. Hier nichts ändern!
\newsavebox{\inhaltkl}
\def\rahmenkl{\sbox{\inhaltkl}\bgroup\small\renewcommand{\baselinestretch}{1}\vbox\bgroup\hsize\textwidth}
\def\endrahmenkl{\par\vskip-\lastskip\egroup\egroup\fboxsep3mm%
\framebox[\textwidth][l]{\usebox{\inhaltkl}}}

% Quelltextrahmen, normale Groesse. Hier nichts ändern!
\newsavebox{\inhalt}
\def\rahmen{\sbox{\inhalt}\bgroup\renewcommand{\baselinestretch}{1}\vbox\bgroup\hsize\textwidth}
\def\endrahmen{\par\vskip-\lastskip\egroup\egroup\fboxsep3mm%
\framebox[\textwidth][l]{\usebox{\inhalt}}}

% Trennvorschläge für falsch getrennte Wörter.
% Wird häufig bei eingedeutschen Wörtern benötigt, da LaTeX hierbei
% gerne falsch trennt. Alternativ kann auch im Fliesstext ein
% Trennvorschlag per "\-" hinterlegt werden, bspw.:
% Die Hard\-ware besteht aus A und B.
\hyphenation{
Hard-ware
}

% Sonstige Befehlsdefinitionen hier ablegen.
\newcommand{\entspricht}{\stackrel{\wedge}{=}}

% Tabellenspaltendefinitionen mit fester Breite --> somit Zeilenumbruch innerhalb einer Zelle möglich
% aus http://www.torsten-schuetze.de/tex/tabsatz-2004.pdf
\usepackage{array, booktabs}
\newcolumntype{f}{>{$}l<{$}}
\newcolumntype{n}{>{\raggedright}l}
\newcolumntype{N}{>{\scriptsize}l}
\newcolumntype{v}[1]{>{\raggedright\hspace{0pt}}m{#1}}
\newcolumntype{V}[1]{>{\scriptsize\raggedright\hspace{0pt}}m{#1}}
\newcolumntype{Z}[1]{>{\raggedright\centering}m{#1}}
\newcolumntype{k}[1]{>{\raggedright}p{#1}}
% ergibt Tabllenspalte fester Breite, linksbündig
% Umbruch innerhalb der Zelle mit \\, neue Tabellezeile mit \tabularnewline
% \addlinespace für Gruppentrennung (aus \texttt{booktabs.sty})


\begin{document}
\onehalfspacing
\begin{titlepage}
\centering
\includegraphics[scale=0.5]{DSB.png}\\[3ex]
{\Large \textsc{Deutscher Schachbund}}\\[3ex]
{\Large Thüringer Schachbund}\\[3ex]
\vfill
{\Large \textbf{\artderausarbeitung}}\\[4ex]
{\large \textbf{\themaderarbeit}}\\[5ex]
\vfill
\begin{tabular}{rl}
\hline \\
vorgelegt von:          & \quad \namedesautors\\[1,5ex]
eingereicht am:         & \quad 10.\,11.\,2020\\[1,5ex]
geboren am:             & \quad 28.\,04.\,1998 in Sondershausen\\[1,5ex]
Ausbildung:             & \quad C-Trainer Lehrgang 2020\\[1,5ex]
Hausarbeit:             & \quad Lehrprobe\\[1,5ex]
                        & \quad Starke und schwache Felder\\[1,5ex]
Dozent:                 & \quad Norbert Reichel (B-Trainer)\\[1,5ex]

\end{tabular}
\vfill
\end{titlepage}








%% ++++++++++++++++++++++++++++++++++++++++++++++++++++++++++++
%% Zusammenfassung, Abstract
%% ++++++++++++++++++++++++++++++++++++++++++++++++++++++++++++


\renewcommand{\abstractname}{Kurzfassung}
\begin{abstract}
\ldots \emph{Hier später die eigene deutsche Kurzfassung einfügen}\ldots
\par{}
Dieses Dokument soll als Gerüst für eigene \LaTeX\ Dokumente dienen und
gleichzeitig Beispiele für häufig verwendete Konstrukte wie Tabellen,
Formeln oder Grafiken liefern. Es empfielt sich, diese Elemente per
Cut\&Paste zu kopieren und einzufügen.
\end{abstract}




% Inhaltsverzeichnis
\cleardoublepage % Seitenumbruch erzwingen vor Änderung des Nummerierungsstils
\pagenumbering{roman} % Nummerierung der Seiten ab hier: i, ii, iii, iv...
\pagestyle{scrheadings} % Ab hier mit Kopf- und Fusszeile
\tableofcontents

% Die einzelnen Kapitel
\cleardoublepage % Seitenumbruch erzwingen vor Änderung des Nummerierungsstils
\pagenumbering{arabic} % Nummerierung der Seiten ab hier: 1, 2, 3, 4...





%%% Inhalt %%%%

\chapter{Training}
\chapter{Schwache Felder}
\chapter{Trainingsplan}
\chapter{Partien}
\chapter{Aufgaben}
\chapter{Lösungen}



%%% Ende %%%

\appendix
\part*{Anhang}




% Literaturverzeichnis einbinden
%% ++++++++++++++++++++++++++++++++++++++++++++++++++++++++++++
%% Anhang: Literaturverzeichnis
%% ++++++++++++++++++++++++++++++++++++++++++++++++++++++++++++


% Mit dem Befehl \nocite werden auch nicht im Text zitierte
% aus der Literaturdatenbank mit in das Literaturverzeichnis aufgenommen.
% Ein "\nocite{*}" übernimmt ungeprüft die komplette Datenbank.
%\nocite{*}

\cleardoublepage
\ihead[]{Literaturverzeichnis}
\bibliographystyle{alphadin}
\bibliography{literatur} % "literatur.bib" ist hier die einzige Literaturdatenbank.

% Alternativ: Mehrere Datenbanken verwenden, falls eine
% oder mehrere umfangreiche Sammlungen exisitieren:
%\bibliography{literatur_buecher,literatur_weblinks}


% Abbildungsverzeichnis einbinden
%% ++++++++++++++++++++++++++++++++++++++++++++++++++++++++++++
%% Anhang: Abbildungsverzeichnis
%% ++++++++++++++++++++++++++++++++++++++++++++++++++++++++++++


% Keine Änderungen vornehmen!
\cleardoublepage
\ihead[]{Abbildungsverzeichnis}
\listoffigures


% Tabellenverzeichnis einbinden
%% ++++++++++++++++++++++++++++++++++++++++++++++++++++++++++++
%% Anhang: Tabellenverzeichnis
%% ++++++++++++++++++++++++++++++++++++++++++++++++++++++++++++

% Hier keine weiteren Änderungen vornehmen
\cleardoublepage
\ihead[]{Tabellenverzeichnis}
\listoftables



% Abkürzungsverzeichnis einbinden
%% ++++++++++++++++++++++++++++++++++++++++++++++++++++++++++++
%% Anhang: Abkürzungsverzeichnis
%% ++++++++++++++++++++++++++++++++++++++++++++++++++++++++++++


% Hier keine weiteren Änderungen vornehmen
\cleardoublepage
\ihead[]{Abkürzungsverzeichnis und Formelzeichen}
\printnomenclature





% Abschlusserklärung
%% ++++++++++++++++++++++++++++++++++++++++++++++++++++++++++++
%% Vorwort: Abschlusserklärung
%% ++++++++++++++++++++++++++++++++++++++++++++++++++++++++++++


\chapter*{Erklärung}
\addcontentsline{toc}{chapter}{Erklärung}
\ihead[]{Erklärung}

Die vorliegende Arbeit habe ich selbstständig ohne Benutzung anderer als der
angegebenen Quellen angefertigt. Alle Stellen, die wörtlich oder sinngemäß
aus veröffentlichten Quellen entnommen wurden, sind als solche
kenntlich gemacht. Die Arbeit ist in gleicher oder ähnlicher Form oder
auszugsweise im Rahmen einer oder anderer Prüfungen noch nicht vorgelegt
worden.
\\[2cm]
Ilmenau, den \today \hfill \namedesautors


\end{document}

